\documentclass[letterpaper]{report}

\usepackage{tikz, background, titling, setspace}
\usepackage{tabu}
\usetikzlibrary{calc}
\usepackage{fancyhdr}
\usepackage{tcolorbox}
\usepackage{layout}
\usepackage[left=2.5cm,right=1cm,top=1cm]{geometry}
\usepackage{pgf}
\usepackage{pgfpages}
\usepackage{pgfplots}
\usepackage{graphicx}
\usepackage{amsfonts}
\usepackage{amsmath}
\usepackage{amssymb}
\usepackage{tikz,pgfplots}
\usepackage[abspage,user,lastpage]{zref}
\usepackage{lastpage}
\usepackage{marginnote}
\usepackage{eso-pic}
\usepackage{ulem}
\usepackage{tikz}
\usepackage{circuitikz}
\usepackage{tikz-timing}
\usepackage{wrapfig}
\usepackage[export]{adjustbox}
\usepackage{listings}
\usepackage{xcolor}
\usepackage{soul}
\usepackage{cite}
\usepackage{url}
\usepackage[absolute]{textpos}
\usepackage{amssymb}  % provides Rbb, etc.
\usepackage{latexsym}  % provides \Box
\usepackage{array}
\usepackage{siunitx}
\usepackage{steinmetz}


% Defined Colors to be used throughout
% Change the color by uncommenting the color you want and commenting out the one you don't
  % Page Background color
    % #FFEC8B (LightGoldenrod1)
        % \definecolor{PageColorz}{RGB}{255,236,139}
    % #FFFE4  (darker YELLOW) -Defualt-
      %\definecolor{PageColorz}{RGB}{255,255,200}
    % #FFFFDC (Light Yellow1) -Recomendable-
      \definecolor{PageColorz}{RGB}{255,255,220}
    % #FFFFED (Light Yellow)
      % \definecolor{PageColorz}{RGB}{255,255,237}
    % #fffff5 (Lighter Yellow)
      % \definecolor{PageColorz}{RGB}{255,255,245}
    % #ffffff (White)
      % \definecolor{PageColorz}{RGB}{255,255,255}
    % #BFEFFF (LightBlue1)
      % \definecolor{PageColorz}{RGB}{191,239,255}
    % #E0FFFF (LightCyan)
      % \definecolor{PageColorz}{RGB}{224,255,255}
    % #104E8B (DodgerBlue4)
	    % \definecolor{PageColorz}{RGB}{16,78,139}
    % #097969 (CadmiumGreen)
      % \definecolor{PageColorz}{HTML}{097969}
    % #AFE1AF (CeladonGreen) 
      % \definecolor{PageColorz}{HTML}{AFE1AF}

    % MEME colorz (PINK) #FFb6C1
      % \defineColor{PageColorz}{HTML}{FFB6C1}
	\pagecolor{PageColorz}

  % Line Colors 
    % #b94249 (Green) -Default-
      \definecolor{MyLines}{RGB}{73,185,66}
    % #6495ed (Blue)
      % \definecolor{MyLines}{RGB}{97,145,179}
    % #104E8B (DodgerBlue4)
      % \definecolor{MyLines}{RGB}{16,78,139}
    % #E0FFFF (LightCyan)
      % \definecolor{MyLines}{RGB}{224,255,255}
    % #000000 (Black)
      % \definecolor{MyLines}{RGB}{0,0,0}
    % #FFFFFF (White)
      % \definecolor{MyLines}{RGB}{255,255,255}

  % Default Text Color
    % #b94249 (Green)
      % \definecolor{MyColor}{RGB}{73,185,66}
    % #6495ed (Blue)
      % \definecolor{MyColor}{RGB}{97,145,179}
    % #000000 (Black) -Default-
      \definecolor{MyColor}{RGB}{0,0,0}
  
    \color{MyColor}

  % GridColor
    % #DFFFFF (UltraLightGrey) -Defualt-
      \definecolor{GridColor}{HTML}{DFFFFF}
    % #D3D3D3 (DarkishGrey)
      %\definecolor{GridColor}{HTML}{D3D3D3}

  % Extra Defined Color
  \definecolor{Extra}{RGB}{16,78,160}
  % \definecolor{Extra}{RGB}{205,245,255}

% A compatability option for pgfplots
\pgfplotsset{compat=1.18}

% Document Properties
\newcommand{\course}{ECE 3620}
\newcommand{\studentName}{}
\newcommand{\aNumber}{A02346768}
\newcommand{\assnNumber}{7}
\newcommand{\w}{\omega}
\newcommand{\Ohm}{\Omega}
\newcommand{\D}{^\circ}
% Draw the Background lines
\backgroundsetup{%
 position=current page.center,
 angle=0,
 scale=1,
 contents={%
  \begin{tikzpicture}%
    [
      normal lines/.style={MyLines, thick},
      grid lines/.style={GridColor, very thin},
      every node/.append style={black, align=center, opacity=1}
    ]
    \foreach \y in {0.5,1,...,27} {\draw[grid lines] (2.5cm,\y cm) -- (20.5cm,\y cm); }
    \foreach \x in {2.5,3,...,20.5}          {\draw[grid lines] (\x cm,0.5cm) -- (\x cm,27 cm); }
	% Draw the background lines of the document
	\draw[normal lines] (3.166in,27cm) -- (3.166in,11in);
	\draw[normal lines] (5.333in,27cm) -- (5.333in,11in);
    \draw[normal lines] (2.5cm,0in) -- (2.5cm,11in);
    \draw[normal lines] (20.5cm,0in) -- (20.5cm,11in);
    \draw[normal lines] (0cm,27cm) -- (21.5cm,27cm);
	% Make Nodes for the header to appear
	\node (t) [font=\footnotesize, anchor=center] at (4.25in, 27.5cm) {\course, Homework: \assnNumber};
    \node (t) [font=\footnotesize, anchor=center] at (5cm, 27.5cm) {October\hspace{2em},2025};
    \node (t) [font=\footnotesize, anchor=center] at (11cm, 27.5cm) {\studentName};
	\node (t) [font=\footnotesize, anchor=center] at (21cm, 27.5cm) {\thepage/\pageref{LastPage}};
  \end{tikzpicture}%
}}


% Pagestyle for no header 
\pagestyle{empty}

% Spacing Controls
\setlength{\parindent}{0pt}

% New environments

\newcounter{ProblemNumber}[section]
\newenvironment{problem}[1][]{
    \noindent
	\textbf{Problem #1 \refstepcounter{ProblemNumber}:} \\ 
    \setlength{\leftskip}{0.125in}
    \setlength{\rightskip}{0.125in}
}
{
	\newline 
	\textcolor{MyLines}{\hrule}
    \vspace{6pt}
    \par
}
%\newcounter{SolutionNumber}[section]
\newenvironment{answer}[1][]{
  \setlength{\leftskip}{0.125in}
  \setlength{\rightskip}{0.125in}
  %\textbf{Solution #1 \refstepcounter{ProblemNumber}:} \\ 
}
{
  \newline
  \textcolor{MyLines}{\hrule}
  \vspace{2pt}
  \textcolor{MyLines}{\hrule}
  \vspace{6pt}
}


\begin{document}

\begin{problem}[{3.30}]
	\\ \textbf{Given:} The values $R_1=1\Omega, R_2=2\Omega,L_1=1mH,$ and $L_2 = 2mH$; and given the circuit below
  \begin{center}
    \includegraphics[scale=0.35]{./pr3.30.png}
  \end{center}
	\textbf{Find:} 
  \begin{itemize}
    \item[a.] $H(s) = \frac{V_o}{V_i}$
    \item[b.] $h(t)$
  \end{itemize}~
\end{problem}
\newpage

\begin{problem}[{3.31}]
	\\ \textbf{Given:} The values $R_1 = 1\Omega$, $R_2= 2\Omega$, $C_1=1\mu F$, and $C_2=2\mu F$; and given the circuit below
  \begin{center}
    \includegraphics[scale=0.35]{./pr3.31.png}
  \end{center}
	\textbf{Find:}
  \begin{itemize}
    \item[a.] $H(s) = \frac{V_o}{V_i}$
    \item[b.] $h(t)$
  \end{itemize}~
\end{problem}
\newpage

\begin{problem}[{3.32}]
	\\ \textbf{Given:} The values $R = 5\Omega$, $L=0.1mH$ , and $C=1\mu F$; and given the circuit below
   \begin{center}
    \includegraphics[scale=0.35]{./pr3.32.png}
  \end{center}
 \begin{itemize}
    \item[a.] $H(s) = \frac{V_o}{V_i}$
    \item[b.] $h(t)$
  \end{itemize}~
  \textbf{Find:}
\end{problem}
\newpage

\begin{problem}[{3.33}]
	\\ \textbf{Given:} An LTI system is described by the LCCDE
  \begin{align}
    \frac{d^2y}{dt^2} - 5 \frac{dy}{dt} + 6y = \frac{dx}{dt}+5x.
  \end{align}
	\textbf{Find:} Is the system BIBO stable?
\end{problem}
\newpage

\begin{problem}[{3.36}]
	\\ \textbf{Given:} An LTI system has transfer function 
  \begin{align}
    H(s) = \frac{(s+1)(s+2)}{s(s+3)}
  \end{align}
	\textbf{Find:} Is it BIBO stable?
\end{problem}
\newpage

\begin{problem}[{4.8}]
  \\ \textbf{Given:} Before closing the switch $v_C(0^-) = 24 V$. Also, the element values are $R=1 \Omega$, $L=0.8 H$, 
  and $C = 0.25 F$.
  \begin{center}
    \includegraphics[scale=0.3]{./pr4.8.png}
  \end{center}
  \textbf{Find:} The equation for $i_L(t)$
\end{problem}
\newpage

\begin{problem}[{4.9}]
	\\ \textbf{Given:}\\
  \begin{center}
    \includegraphics[scale=0.3]{./pr4.9.png}
  \end{center}

	\textbf{Find:}
\end{problem}
\newpage

\begin{problem}[{4.11}]
	\\ \textbf{Given:}\\
  \begin{center}
    \includegraphics[scale=0.3]{./pr4.11.png}
  \end{center}

	\textbf{Find:}
\end{problem}
\newpage

\begin{problem}[{4.32}]
	\\ \textbf{Given:}\\
  \begin{center}
    \includegraphics[scale=0.3]{./pr4.32.png}
  \end{center}

	\textbf{Find:}
\end{problem}
\newpage

\begin{problem}[{11.}]
	\\ \textbf{Given:}\\
  \begin{center}
    \includegraphics[scale=0.3]{./pr11.png}
  \end{center}
	\textbf{Find:}
\end{problem}

\end{document}
